\subsubsection{Geometrische Merkmalsextraktion}

Die Extraktion geometrischer Merkmale basiert sowohl auf den 2D-Kronenpolygone als auf den 3D-Messpunkten. Als Bindeglied zwischen Vektordaten und Punktwolken-Segmenten dient dabei die eindeutige ID \cite{munzingerMappingUrbanForest2022}. Zunächst wird durch eine direkte geometrische Auswertung der Polygone die Kronenfläche berechnet \cite{yuIndividualTreeSegmentation2024}. Ausgehend davon lässt sich unter der Annahme einer idealisierten Kreisform anschließend der Kronenradius ableiten.

Die vertikalen und volumenbasierten  Strukturparameter erfordern hingegen eine detaillierte Analyse der zugeordneten 3D-Punkte. Zur Bestimmung der Kronenansatzhöhe wird ein histogrammbasierter Ansatz gewählt, da dieser robust gegenüber Rauschen sowie der abschattungsbedingten geringen Punktdichte im Stammbereich ist. Hierzu teilt der Algorithmus die Höhewerte der Punkte eines Baumes in vertikale Intervalle mit einer Schichtdicke von 0,3 m ein und analysiert die Punktdichte von unten nach oben. Die Grenze des Kronenansatzes wird an der Schwelle definiert, an der die Punktdichte erstmals mindestens zehn Prozent des Maximalwerts innerhalb des Baums erreicht. Dieser Schwellwert markiert somit den strukturellen Übergang vom spärlich erfassten Stammbereich zum dichten Kronenraum \cite{zotero-item-73}. Die Kronenhöhe ergibt sich folglich aus der Differenz zwischen der absoluten Höhe der Baumspitze und der ermittelten Kronenansatzhöhe. Das Kronenvolumen wird berechnet, indem eine konvexe Hülle um alle Punkte der jeweiligen Baumkrone gespannt wird \cite{munzingerMappingUrbanForest2022}.

Abschließend werden alle extrahierten Parameter über die Baum-ID als Attribute in den Layer der GPKG-Datei, welcher bereits die Baumspitzen enthält, überführt und der Datensatz somit zeilenweise angereichert.